\documentclass[tikz,border=3pt]{standalone}
\begin{document}

\begin{tikzpicture}

% 軸
\draw[->] (0,0) -- (0,7);
\draw[->] (0,0) -- (7,0);

% 枠（x=1, y=1）
\draw (0,6) -- (6,6) -- (6,0);

% 軸ラベル（←ここをLargeに）
\node[font=\Large] at (7,-0.5) {X};
\node[font=\Large] at (-0.5,7) {Y};

% --- x軸の目盛り ---
\node[font=\Large] at (0,-0.5) {$0$};
\node[font=\Large] at (36/7,-0.5) {$\frac{6}{7}$};
\node[font=\Large] at (6,-0.5) {$1$};

% --- y軸の目盛り ---
\node[font=\Large] at (-0.5,0) {$0$};
\node[font=\Large] at (-0.5,18/7) {$\frac{3}{7}$};
\node[font=\Large] at (-0.5,36/7) {$\frac{6}{7}$};
\node[font=\Large] at (-0.5,6.2) {$1$};

% 対角線 y = x
\draw (0,0) -- (36/7,36/7);
\draw[dashed] (36/7,36/7) -- (6,6);

% --- 水平線 ---
% y = 3/7
\draw (0,18/7) -- (36/7,18/7);
\draw[dashed] (36/7,18/7) -- (6,18/7);

% y = 6/7
\draw (0,36/7) -- (36/7,36/7);
\draw[dashed] (36/7,36/7) -- (6,36/7);

% --- 縦線 x = 6/7（y <= 6/7 のみ点線）---
\draw[dashed] (36/7,0) -- (36/7,36/7);

% 曲線1: y = (7/12)x^2 （x <= 6/7）
\draw[domain=0:6/7, smooth, variable=\x]
  plot ({6*\x}, {6*(7/12)*\x*\x});

% 曲線2: y = 1 - 7/6 (1/2 x^2 - 7/18 x^3) （x <= 6/7）
\draw[domain=0:6/7, smooth, variable=\x]
  plot ({6*\x}, {6*(1 - (7/6)*(1/2*\x*\x - 7/18*\x*\x*\x))});

\end{tikzpicture}

\end{document}